\documentclass[10pt,a4paper]{amsart}
\usepackage[margin=19mm]{geometry}
\usepackage{amsmath,amssymb,mathtools}
\usepackage[scr=rsfs,cal=euler]{mathalfa}
\usepackage{xfrac}
\usepackage[unicode,hidelinks,bookmarks=false]{hyperref}
\newcommand{\ie}{i.e.\ }
\newcommand{\cf}{cf.\ }
\newcommand{\resp}{respectively\ }
\newcommand{\Subsection}{\subsection}
\providecommand{\nobreakdash}{\nobreak\hbox{-}}
\setlength{\emergencystretch}{2em}
\multlinegap0pt
\usepackage{float}
\usepackage{amssymb}
\usepackage{enumerate}
\usepackage{bm}
\newtheorem{thm}{Theorem}
\newtheorem{claim}{Claim}
\newcommand{\C}{\mathbb{C}}          
\newcommand{\Cal}{\mathcal}
\renewcommand{\Im}{\mathrm{Im}}
\newcommand{\Mp}{M^{\fg}_{\fp}}
\newcommand{\Mpp}{M^{\fg'}_{\fp'}}
\newcommand{\Nn}{N_n^-}
\newcommand{\Np}{(N^-)'}
\newcommand{\Z}{\mathbb{Z}}
\newcommand{\fg}{{\mathfrak g}}
\newcommand{\fn}{{\mathfrak n}}
\newcommand{\fp}{{\mathfrak p}}
\title{Differential symmetry breaking operators from a line bundle to a vector bundle over real projective spaces}
\author{Toshihisa Kubo}
\date{Source: arXiv:2408.01213v3 (CC BY 4.0)}
\begin{document}
\maketitle

\noindent\emph{Source and licence.} Differential symmetry breaking operators from a line bundle to a vector bundle over real projective spaces. Version arXiv:2408.01213v3 (CC BY 4.0), distributed by arXiv under the Creative Commons Attribution 4.0 International licence, http://creativecommons.org/licenses/by/4.0/, as recorded in the arXiv licence metadata for this exact version and shown on the abstract page https://arxiv.org/abs/2408.01213v3. Full licence notice: \texttt{licenses/arxiv-2408.01213-CC-BY-4.0.txt}.\par\smallskip

\noindent\textit{Editorial scope.} Counted unit: the article's thm thm:GVM31b (source0.tex lines 4812--4824). The article's complete original proof of it is delivered with it (source0.tex lines 4826--4955, one \texttt{proof} environment of 130 lines). No mathematics is written, repaired or re-derived: the statement and the proof are the article's own lines, in the article's own notation, and no retained line is substituted or re-flowed.  Retained uncounted as the source-local context the proof needs: lines 4632--4646; lines 4733--4737; lines 4739--4743; lines 4792--4800; lines 4802--4810.  Nothing was omitted from any retained span, so the package carries no omission record.  No retained line contains a citation, so the wrapper carries no bibliography and no bibliography entry.  No retained line contains a figure environment, an image-inclusion command or drawing code, so no image is delivered and the package declares no asset dependency.  Editorial formatting only, in addition: an AMS \texttt{amsart} preamble that restates the article's own declarations and macros used by the retained lines, namely the environments \texttt{thm}, \texttt{claim} on separate counters, together with the standard packages those lines need.  The retained environments print in this document as Theorem 1, Claim 1 in order of appearance, whereas the article prints the same items with the numbering of its own full document.  The complete native source of the article as distributed in the arXiv e-print for this exact version is bundled unchanged as \texttt{sources/163264/source0.tex}.
\begin{thm}\label{thm:GVM2}
Let $n\geq 2$. For $p \in \Z_{\geq 0}$, 
the following isomorphism holds
in the Grothendieck group of $\Cal{O}^{\fp'}$:
\begin{equation}\label{eqn:GVM3a}
[\Im(\varphi_{p+1})\vert_{\fg'}] \simeq 
\bigoplus_{d=0}^p [\Im (\varphi_{d+1}')] \oplus \bigoplus_{j \geq 1} [\Mpp(-j)].
\end{equation}
Further, for $n=2$, we have 
\begin{equation}\label{eqn:GVM3b}
[\Im(\varphi_{p+1})\vert_{\fg'}] \simeq 
\bigoplus_{d=0}^p 2 \cdot [\Mpp(-(d+2))] \oplus [\Mpp(-1)] \oplus 
\bigoplus_{j\geq p+3} [\Mpp(-j)].
\end{equation}
\end{thm}

\begin{equation}\label{eqn:M300}
\C^m[N_n^-] \otimes \C_s \subset \Mp(s)^{\fn_+'} 
\quad
\text{for all $m \in \Z_{\geq 0}$}.
\end{equation}

\begin{equation}\label{eqn:M30a}
\C^m[N_n^-] \otimes \C_s\simeq \C_{s-m}
\quad
\text{for all $m \in \Z_{\geq 0}$},
\end{equation}

\begin{align}
\Im(\varphi_{p+1}) &= 
\Cal{U}(\fg)\left(
\C^{p+1}[(N^-)', N_n^-] \otimes \C_p
\right), \label{eqn:M30A}\\
\Im(\varphi'_{d+1}) &= \Cal{U}(\fg')\left(
\C^{d+1}[(N^-)'] \otimes \C_d
\right),\label{eqn:M30B}
\end{align}

\begin{equation}\label{eqn:M30D}
\C^{p+1}[(N^-)', N_n^-] \otimes \C_p
\subset \Mp(p)^{\fn_+}
\quad
\text{and}
\quad
\C^{d+1}[(N^-)'] \otimes \C_d
\subset \Mpp(d)^{\fn_+'}.
\end{equation}

\begin{thm}\label{thm:GVM31b}
Let $n\geq 2$. For $p \in \Z_{\geq 0}$, we have 
\begin{equation}\label{eqn:GVM31A}
\Im(\varphi_{p+1})\vert_{\fg'} \simeq 
\bigoplus_{d=0}^p \Im (\varphi_{d+1}') \oplus \bigoplus_{j \geq 1} \Mpp(-j).
\end{equation}
In partucular, for $n=2$, we have 
\begin{equation}\label{eqn:GVM31B}
\Im(\varphi_{p+1})\vert_{\fg'} \simeq 
\bigoplus_{d=0}^p 2 \cdot \Mpp(-(d+2)) \oplus \Mpp(-1) \oplus 
\bigoplus_{j\geq p+3} \Mpp(-j).
\end{equation}
\end{thm}

\begin{proof}
Since \eqref{eqn:GVM31B} follows from \eqref{eqn:GVM31A}
as in Theorem \ref{thm:GVM2}, it suffices to show \eqref{eqn:GVM31A}.

First observe that
\begin{align}
\C^{p+1}[(N^-)', N_n^-]\otimes \C_p
&=\sum_{b=0}^{p+1}\C^b[\Np]\C^{p+1-b}[\Nn]\otimes \C_p \nonumber \\
&\simeq \sum_{b=0}^{p+1}\C^b[\Np]\otimes \C_{b-1} \nonumber \\
&=\sum_{d=0}^p \C^{d+1}[\Np] \otimes \C_d + (1\otimes \C_{-1}).
\label{eqn:M30C}
\end{align}
We note that \eqref{eqn:M30a} is applied from line one to line two.
Then, by \eqref{eqn:M30A} and \eqref{eqn:M30C}, we have
\begin{align}
\Im(\varphi_{p+1})
&=\Cal{U}(\fg)\left(\C^{p+1}[(N^-)', N_n^-] \otimes \C_p\right) \nonumber \\
&=\C[\Np,\Nn]\C^{p+1}[(N^-)', N_n^-] \otimes \C_p \nonumber \\
&=\sum_{d=0}^p \C[\Np,\Nn]\C^{d+1}[\Np] \otimes \C_d 
+ \C[\Np,\Nn]\otimes \C_{-1} \label{eqn:M30E}\\
&=: \text{(B)} + \text{(A)}.\label{eqn:M31A}
\end{align}

The second term \text{(A)} is given as
\begin{align}
\text{(A)}
= \C[\Np,\Nn]\otimes \C_{-1}
&=\bigoplus_{c=0}^\infty \C[\Np]\C^c[\Nn]\otimes \C_{-1} \nonumber\\
&\simeq \bigoplus_{c=0}^\infty\C[\Np]\otimes \C_{-1-c} \nonumber\\
&=\bigoplus_{j=1}^\infty\C[\Np]\otimes \C_{-j} \nonumber\\
&=\bigoplus_{j=1}^\infty\Mpp(-j).\label{eqn:M30F}
\end{align}
By \eqref{eqn:M300}, we have
$\C_{-j}\simeq \C^c[\Nn]\otimes \C_{-1} \subset \Mp(-1)^{\fn_+'}$,
which verifies the identity from line three to line four.

Next, for \text{(B)}, we have 
\begin{align}
\text{(B)}
&=\sum_{d=0}^p \C[\Np,\Nn]\C^{d+1}[\Np] \otimes \C_d \nonumber \\
&=\sum_{d=0}^p\sum_{c=0}^\infty \C[\Np]\C^c[\Nn]\C^{d+1}[\Np]\otimes \C_d
\nonumber\\
&=\sum_{d=0}^p\C[\Np]\C^{d+1}[\Np]\otimes \C_d
+
\sum_{d=0}^p\sum_{c=1}^\infty \C[\Np]\C^c[\Nn]\C^{d+1}[\Np]\otimes \C_d
\nonumber \\
&=:\text{(B1)} + \text{(B2)}.\label{eqn:M30G}
\end{align}

\noindent
By \eqref{eqn:M30B} and \eqref{eqn:M30D},  \text{(B1)} is given by
\begin{align}
\text{(B1)}
=\sum_{d=0}^p\C[\Np]\C^{d+1}[\Np]\otimes \C_d
&=\bigoplus_{d=0}^p\Cal{U}(\fg')\left(\C^{d+1}[\Np]\otimes \C_d\right)\nonumber \\
&=\bigoplus_{d=0}^p \Im(\varphi'_{d+1}).\label{eqn:M30H}
\end{align}

The actual realizations of $\Mpp(-j)$ in \text{(A)} and $\Im(\varphi'_{d+1})$ in \text{(B1)}
imply that 
the sum $\text{(B1)} + \text{(A)}$ is indeed a direct sum
$\text{(B1)} \oplus \text{(A)}$.
We then claim the following.

\begin{claim}\label{claim:M30}
We have $\text{(B2)} \subset \text{(A)} \oplus \text{(B1)}$.
\end{claim}

Suppose that Claim \ref{claim:M30} holds. Then, 
by \eqref{eqn:M31A}, \eqref{eqn:M30F}, \eqref{eqn:M30G}, and \eqref{eqn:M30H},
we have 
\begin{align*}
\Im(\varphi_{p+1}) 
= \text{(B)}+\text{(A)}
&=\text{(B1)} + \text{(B2)} + \text{(A)}\\
&= \text{(B1)}\oplus \text{(A)}\\
&\simeq 
\bigoplus_{d=0}^p \Im(\varphi'_{d+1}) \oplus \bigoplus_{j=1}^\infty\Mpp(-j),
\end{align*}
which is what we wish to show. Thus, in the rest of the proof, we aim to show 
Claim \ref{claim:M30}.

Write
\begin{equation*}
\text{(C)}=\C[\Np]\C^c[\Nn]\C^{d+1}[\Np]\otimes \C_d,
\end{equation*}
so that 
$\text{(B2)} = \sum_{d=0}^p\sum_{c=1}^\infty \text{(C)}$.

First observe that \text{(A)}, \text{(B)}, and \text{(C)} are given by
\begin{align*}
\text{(A)}
&= \C[\Np,\Nn]\otimes \C_{-1}\\
&=\C[\Np,\Nn]\C^{p+1}[\Nn]\otimes \C_p,\\[7pt]
\text{(B1)}
&=\sum_{d=0}^p\C[\Np]\C^{d+1}[\Np]\otimes \C_d\\
&=\sum_{b=1}^{p+1}\C[\Np]\C^{b}[\Np]\C^{p+1-b}[\Nn]\otimes \C_p, \\[7pt]
\text{(C)}
&= \C[\Np]\C^c[\Nn]\C^{d+1}[\Np]\otimes \C_d\\
&= \C[\Np]\C^c[\Nn]\C^{b}[\Np]\C^{p+1-b}[\Nn]\otimes \C_p\\
&= \C[\Np]\C^{b}[\Np]\C^{p+1+c-b}[\Nn]\otimes \C_p.
\end{align*}

\noindent
Thus, if $c-b\geq 0$, then
\begin{align*}
\text{(C)}
&=\C[\Np]\C^{b}[\Np]\C^{p+1+c-b}[\Nn]\otimes \C_p\\
&\subset\C[\Np,\Nn]\C^{p+1}[\Nn]\otimes \C_p\\
&=\text{(A)}.
\end{align*}

\noindent
If $c-b < 0$, then $-(p+1) \leq c-b \leq -1$. Thus,
the number $p+1+c-b$ is of the form
\begin{equation*}
p+1+c-b = p+1-b_0
\end{equation*}
for some $b_0 \in \{1,\ldots, p+1\}$. Therefore,
\begin{align*}
\text{(C)}
&=\C[\Np]\C^{b}[\Np]\C^{p+1-b_0}[\Nn]\otimes \C_p\\
&\subset 
\sum_{b=1}^{p+1}\C[\Np]\C^{b}[\Np]\C^{p+1-b}[\Nn]\otimes \C_p\\
&=\text{(B1)}.
\end{align*}

\noindent
As $\text{(B2)} = \sum_{d=0}^p\sum_{c=1}^\infty \text{(C)}$, this shows the claim.
\end{proof}

\end{document}
